\documentclass[12pt]{article}

\usepackage[margin=1in]{geometry}
\usepackage{setspace}
\usepackage{graphicx}
\usepackage{booktabs}
\usepackage{tabularx}
\usepackage{array}
\usepackage{amsmath}
\usepackage{hyperref}
\usepackage{float}
\usepackage{enumitem}

\hypersetup{hidelinks}
\singlespacing
\setlength{\parindent}{0pt}
\setlength{\parskip}{0.5em}
\setlength{\intextsep}{8pt}
\setlength{\textfloatsep}{10pt}
\sloppy

\title{STATS 571 Final Project Report\\Causal Discovery for Nonstationary Financial Time Series}
\author{Deepan Islam \\ Urvi Mehta \\ Santosh Desai}
\date{\today}

\begin{document}
\maketitle

\textbf{Project summary.} Our goal was to reproduce the CD-NOTS paper and then write the final
report using the notation and logic from class. After checking the paper against the code and the
saved outputs, we do not think the current repository is a full scientific replication. It does
recover the broad workflow: lag construction, a context node, graph pruning by conditional
independence tests, orientation rules, and graph artifacts for downstream analysis. But the main
fidelity gap is still important. The advertised KCIT, RCoT, and CMIknn options still route back to
partial correlation, stage 4 is an approximate closure step, and the saved benchmark does not
reproduce the paper's claim that CD-NOTS beats PCMCI. So the honest write-up is a course-oriented
causal analysis built on a reproducible CD-NOTS-style scaffold.

\section*{1. Problem Statement and Causal Question}

The paper studies causal discovery for financial time series that are nonstationary and contain
lagged dependence. That is a natural final-project setting for a causal inference course because it
forces the same questions that show up in observational studies more generally. There is no
randomized treatment assignment. Economically important variables move together. Overlap can be weak
in some parts of the sample. Hidden bias is always plausible.

Our project asks two linked questions. First, does the repository reproduce the paper's discovery
workflow well enough to recover a useful working graph on financial and macroeconomic data? Second,
once we have that working graph, can we translate the downstream analysis into the language of class
by stating an estimand, an identification argument, an adjustment set, and the main threats to that
argument?

As we discussed in Lecture 13 titled: \textit{Recap: Confounding and Causal DAGs}, a learned graph
is only useful if it helps us separate causal paths from noncausal ones. As we discussed in Lecture
15 titled: \textit{Methods for adjustment beyond stratification}, an observational estimate only
earns a causal interpretation after we say what assumptions identify it. That is the standard we
use throughout the report.

\section*{2. Data and Effective Unit of Analysis}

The repository uses two public datasets.

\textbf{Fama-French plus Apple.} The file \texttt{famafrench\_apple\_daily.csv} contains 6,016
daily observations from 2000-01-04 through 2023-11-30. The financial variables used in the case
study are \texttt{Mkt\_RF}, \texttt{SMB}, \texttt{HML}, \texttt{RMW}, \texttt{CMA}, \texttt{RF},
and \texttt{AAPL\_RET}.

\textbf{Multi-country macro panel.} The file \texttt{macro\_countries\_monthly.csv} contains 1,800
monthly observations from 2000-01-31 through 2024-12-31, with 300 rows each for \texttt{CA},
\texttt{FR}, \texttt{GB}, \texttt{IT}, \texttt{JP}, and \texttt{US}. The main variables are
unemployment, CPI, and PPI.

The effective unit is not a person, hospital, or firm. After lagging and preprocessing, the
analysis unit is a date-specific record together with the observed history carried into that row.
That is already a warning sign from a class perspective. The usual superpopulation story is much
less natural here than it is in the standard i.i.d. examples from lecture, so every identifying
assumption has to be read with that limitation in mind.

\section*{3. Case Setup in Class Notation}

As we discussed in Lecture 15 titled: \textit{Methods for adjustment beyond stratification}, the
write-up has to define \(Z\), \(Y\), \(X\), and the assignment space before it talks about
estimators.

\subsection*{3.1 Financial case}

Let \(t=1,\dots,T_F\) index trading days after preprocessing. In the main financial regression, we
set
\[
Z_t = SMB_t, \qquad Y_t = HML_t.
\]
Let \(X_t \in \mathcal{X}_F\) denote the observed covariates and lagged history available at day
\(t\). In this project, that means the observed factor returns, the risk-free rate, the time
context carried by the graph pipeline, and any lagged variables produced upstream. Let
\(L_t \subseteq X_t\) denote the graph-derived adjustment set. In the saved run, the class-oriented
adjustment step uses \(L_t = \{RMW_t\}\).

To keep the notation close to class, let \(\Omega_F\) denote the set of feasible financial exposure
paths \((z_1,\dots,z_{T_F})\) over the study window. Because this is observational time-series
data, \(\Omega_F\) is not known from a design the way it would be in a randomized experiment. It is
only a conceptual assignment space. The class-style estimand is the dose-response contrast
\[
\tau_F(z,z') = E\{Y_t(z) - Y_t(z')\},
\]
where \(Y_t(z)\) is the potential value of \(HML_t\) under an intervention setting \(SMB_t=z\).

The identifying assumptions are the usual ones, but translated carefully to this setting.
Consistency means the observed \(HML_t\) equals \(Y_t(Z_t)\). SUTVA means the potential outcome at
time \(t\) does not depend on unmodeled assignments outside the lagged history already encoded in
the row. Conditional exchangeability means
\[
\{Y_t(z): z \in \Omega_F\} \perp Z_t \mid X_t.
\]
The overlap statement from lecture is harder because the exposure is continuous. The practical
reading is that, within relevant values of \(X_t\), the data must contain enough variation in SMB
that the regression is not pure extrapolation.

\subsection*{3.2 Macro case}

Let \(t=1,\dots,T_M\) index the monthly US observations used in the saved macro run. Here the focal
variables are
\[
Z_t = unemployment_t, \qquad Y_t = \Delta cpi_t,
\]
with the differenced CPI series stored back into the \texttt{cpi} column before estimation. Let
\(X_t \in \mathcal{X}_M\) denote the observed macro history available at month \(t\), including PPI,
time context, and the lagged information retained by preprocessing. In the saved run the graph does
not produce a clean adjustment set near the treatment-outcome pair, so \(L_t\) is effectively empty
and the script raises an ambiguity warning.

Let \(\Omega_M\) denote the set of feasible unemployment paths over the study window. Again this is
conceptual rather than design-based. The natural class-style estimand is
\[
\tau_M(z,z') = E\{Y_t(z) - Y_t(z')\}.
\]
Consistency means the observed CPI increment equals the potential outcome under the realized
unemployment level. SUTVA means that, once the chosen lag structure is fixed, the potential CPI
increment at month \(t\) does not depend on unmodeled spillovers or outside shocks that alter the
assignment process. Conditional exchangeability would require
\[
\{Y_t(z): z \in \Omega_M\} \perp Z_t \mid X_t,
\]
which is a strong claim in macro data. The overlap issue is also sharper here because the sample is
small and trending.

\begin{table}[H]
\centering
\small
\begin{tabularx}{\textwidth}{p{0.18\textwidth} X X}
\toprule
\textbf{Object} & \textbf{Financial case} & \textbf{Macro case} \\
\midrule
Unit & Trading day \(t\) after preprocessing and lagging & US month \(t\) after preprocessing \\
\(Z\) & \(SMB_t\) & \(unemployment_t\) \\
\(Y\) & \(HML_t\) & \(\Delta cpi_t\), stored in the \texttt{cpi} column after differencing \\
\(X\) & Observed factor returns, risk-free rate, time context, and lagged history available at \(t\) & Observed macro covariates, time context, PPI, and lagged history available at \(t\) \\
\(L\) & Graph-derived set; saved run uses \(\{RMW_t\}\) & No clean set in the saved run because the graph stays partially undirected \\
\(\Omega\) & Feasible SMB exposure paths over the sample window & Feasible unemployment paths over the sample window \\
SUTVA-style reading & No unmodeled interference across dates once the chosen lags absorb the relevant history & No unmodeled macro spillovers beyond the chosen monthly history \\
\bottomrule
\end{tabularx}
\caption{The two case studies written in the notation from class.}
\end{table}

\section*{4. Working DAG, Backdoor Paths, and Adjustment}

As we discussed in Lecture 13 titled: \textit{The Backdoor Criterion}, confounding is easiest to
see by thinking about paths. A backdoor path from \(Z\) to \(Y\) is any noncausal path that starts
with an arrow into \(Z\). If conditioning on a set \(L\) blocks every open backdoor path, and \(L\)
contains no descendants of treatment, then \(L\) satisfies the backdoor criterion.

That is exactly how we read the discovery output. The graph returned by the repository is not known
truth. It is an estimated graph. Still, once we have a directed subgraph, we can use it as a
working DAG for choosing candidate controls. The script \texttt{class\_oriented\_graph\_adjustment.py}
does this directly: it reads the directed CD-NOD output, builds a candidate adjustment set from the
parents of \(Z\), and fits
\[
E(Y_t \mid Z_t, L_t) = \beta_0 + \beta_Z Z_t + \beta_L^\top L_t.
\]

The important limitation is the same one we discussed in class. If the graph is wrong, then the
adjustment set can be wrong. If the graph remains partially directed near treatment and outcome,
identification is unsettled. The macro case shows this clearly. The saved run keeps unemployment
and CPI as undirected neighbors, so the script correctly flags the result as graph-ambiguous.

\section*{5. Estimation Strategy and What It Is Not}

As we discussed in Lecture 16 titled: \textit{Recap: Outcome regression / model standardization},
the downstream estimator in this repository is best understood as outcome regression. Under
consistency, conditional exchangeability, and overlap, the identifying formula is
\[
E\{Y_t(z)\} = E_{X_t}\left[E\{Y_t \mid Z_t=z, X_t\}\right].
\]
The implementation uses this logic in a simple linear form with HAC standard errors to account for
serial dependence.

This matters because it tells us what claim we are and are not making. The saved coefficients are
not IPW estimates. They are not AIPW estimates. They are not TMLE estimates. As we discussed in
Lecture 17 titled: \textit{Recap: Outcome Regression and Inverse Probability Weighting} and Lecture
18 titled: \textit{Recap: Augmented Inverse Probability Weighting}, the repository never fits a
propensity score model and never adds the augmentation term that gives AIPW its double robustness.
So there is no insurance against one model being wrong. We have a graph-guided outcome regression,
not a doubly robust observational-study pipeline.

\section*{6. What the Code Actually Implements}

The paper describes CD-NOTS as a four-stage procedure, and the repository follows that outline at a
high level:
\begin{enumerate}[leftmargin=1.5em]
\item build lagged variables and add a context node,
\item prune a dense graph with conditional independence tests,
\item orient edges using time order, lag order, and v-structures,
\item apply a final orientation step.
\end{enumerate}

The fidelity gap shows up in the details. In \texttt{ci\_tests.py}, the labels \texttt{kcit},
\texttt{rcot}, and \texttt{cmiknn} still call the same fallback routine that uses partial
correlation. Stage 4 also differs from the paper. The code uses a Meek-style closure approximation
rather than the paper's change-in-mechanism rule.

There is also a second graph path in \texttt{discovery2}. That workflow runs \texttt{causal-learn}
CD-NOD directly on the datasets and exports directed and undirected edge files. The saved
class-oriented regressions consume those graph artifacts. So the project has two layers: a
CD-NOTS-style implementation for simulation and case-study workflows, and a CD-NOD artifact
pipeline that feeds the downstream adjustment scripts.

\section*{7. Results}

\subsection*{7.1 Quick simulations and benchmark}

The saved simulation files come from a quick run with node grids \(\{3,5\}\), observation grids
\(\{50,150\}\), and five graphs per setting. The F1 scores are identical across
\texttt{parcorr}, \texttt{kcit\_hbe}, \texttt{rcot\_hbe}, and \texttt{cmiknn}: 0.333 for 3 nodes
with 50 observations, 0.200 for 5 nodes with 50 observations, 0.300 for 3 nodes with 150
observations, and 0.408 for 5 nodes with 150 observations. Once we know that the non-ParCorr
labels still fall back to ParCorr, this pattern is exactly what we should expect.

The benchmark against PCMCI is the clearest sign that the replication is incomplete. The file
\texttt{benchmark\_pcmci.csv} reports precision, recall, and F1 of 0 for the current CD-NOTS
implementation and 1 for PCMCI-ParCorr. That result runs against the paper's headline comparison,
so it should be stated plainly rather than explained away.

\begin{figure}[H]
\centering
\includegraphics[width=0.60\textwidth]{results/figures/simulation_f1.png}
\caption{Quick-run simulation F1 output. The near overlap of the curves matches the code audit
because the non-ParCorr labels currently fall back to ParCorr.}
\end{figure}

\subsection*{7.2 Financial case}

The financial case is the strongest part of the project. The case-study edge file reports 95
discovered edges for 2000 to 2007, 77 for 2008 to 2015, 59 for 2016 to 2022, and 66 for the
full-period run. That is at least directionally consistent with the paper's claim that dependence
structure changes over time.

The saved class-oriented adjustment file \texttt{class\_oriented\_adjust\_famafrench\_SMB\_HML.csv}
estimates a coefficient of 0.271 for SMB on HML with standard error 0.034 and
\(p < 1.7 \times 10^{-15}\), using \(L=\{RMW\}\) as the graph-derived adjustment set. A second
file, \texttt{class\_oriented\_adjust\_famafrench\_RMW\_Mkt\_RF.csv}, estimates a coefficient of
\(-0.829\) for RMW on Mkt-RF with \(L=\{HML, SMB\}\).

The caution comes from the diagnostics. The graph falsification report passes only 6 of 13 implied
independence checks. The rolling stability analysis shows the SMB to HML edge in only 8 of 57
windows, while the adjusted coefficient is statistically significant in 31 of 57 windows. So the
regression signal is more stable than the edge itself, which is exactly the kind of tension that a
class-style write-up should acknowledge.

\begin{figure}[H]
\centering
\includegraphics[width=0.55\textwidth]{results/figures/rolling_stability_famafrench_SMB_HML.png}
\caption{Rolling-window stability for the SMB to HML financial analysis. The coefficient appears
more stable than the edge itself.}
\end{figure}

\subsection*{7.3 Macro case}

The macro case is much less decisive. The saved CD-NOD graph directs \texttt{context\_country} into
unemployment, CPI, and PPI, but keeps the pairwise edges among those three variables undirected.
The downstream class-oriented regression still returns a coefficient of \(-0.074\) with standard
error 0.028 and \(p = 0.0076\), but it also raises an ambiguity warning because unemployment and
CPI remain undirected neighbors.

That is exactly how this result should be written up using the lecture language. There is an
adjusted association, but the backdoor argument is not settled. The graph never gives us a clean
\(L\) near the treatment-outcome pair, so the causal interpretation is weaker than in the financial
case.

\begin{figure}[H]
\centering
\includegraphics[width=0.48\textwidth]{results/graphs/case_macro_US.png}
\caption{Saved macro graph for the US case. The remaining undirected edges near unemployment and CPI
are the reason the macro identification argument stays weaker.}
\end{figure}

\begin{table}[H]
\centering
\small
\begin{tabularx}{\textwidth}{p{0.24\textwidth} p{0.22\textwidth} X}
\toprule
\textbf{Analysis} & \textbf{Saved Result} & \textbf{Interpretation} \\
\midrule
Quick simulations & F1 ranges from 0.200 to 0.408, identical across CI labels & Matches the code audit because the non-ParCorr labels fall back to ParCorr \\
PCMCI benchmark & CD-NOTS-ParCorr: \(0/0/0\); PCMCI-ParCorr: \(1/1/1\) & The current benchmark does not support the paper's comparative claim \\
Financial adjustment & SMB on HML coefficient \(= 0.271\), SE \(= 0.034\) & Strong adjusted association under the working graph-derived set \(L=\{RMW\}\) \\
Graph falsification & 6 of 13 implied independence checks pass & Mixed support for the financial working DAG \\
Rolling stability & Edge appears in 8 of 57 windows; coefficient significant in 31 of 57 windows & The coefficient is more stable than the edge \\
Macro adjustment & Unemployment on \(\Delta cpi\) coefficient \(= -0.074\), SE \(= 0.028\) & Association is nonzero, but identification remains ambiguous because the graph is partially directed \\
\bottomrule
\end{tabularx}
\caption{Main numerical outputs used in the discussion.}
\end{table}

\section*{8. Hidden Bias, Overlap, and Robustness}

As we discussed in Lecture 14 titled: \textit{Rosenbaum's model for hidden bias in matched
observational studies}, the right question is not only whether a coefficient is nonzero, but also
how sensitive the conclusion is to unmeasured confounding. The repository does \emph{not} implement
a formal \(\Gamma\)-based sensitivity analysis, so we should not use that language too casually.

What the repository does provide are diagnostics that move in the same direction: graph implication
falsification, rolling-window stability, ambiguity flags when the graph remains partially directed,
and partial-\(R^2\)-style summaries. These are useful checks, but they are not a substitute for a
formal hidden-bias analysis.

Overlap is also more informal than it would be in a lecture-style IPW analysis. The financial case
has many more observations, so the support problem is milder. The macro case has far fewer
observations and stronger trend structure, so the overlap claim is weaker there. That is one reason
the macro coefficient should be treated much more cautiously.

\section*{9. Discussion and Conclusion}

Using the language of class, our final position is simple.

\begin{enumerate}[leftmargin=1.5em]
\item The project has a clear estimand in potential-outcomes notation once we specify \(Z\), \(Y\),
\(X\), and \(\Omega\) for each case.
\item The identification strategy is a backdoor-adjustment argument based on a working DAG, not a
known DAG.
\item The estimator is graph-guided outcome regression with HAC standard errors.
\item The main threats are failures of consistency, SUTVA, exchangeability, overlap, and hidden
bias.
\end{enumerate}

On the software side, the repository is useful and reproducible. On the scientific side, it is not
yet a full replication of the paper because the CI machinery is still simplified, the stage-4
orientation is approximate, and the benchmark does not reproduce the paper's main comparison. So
the right conclusion is narrower than the paper's conclusion. This project gives us a workable
bridge from causal discovery to course-style causal reasoning. It does not yet settle the paper's
empirical claims.

\clearpage
\section*{References}

\footnotesize
\begin{itemize}[leftmargin=1.5em, itemsep=0pt, parsep=0pt, topsep=2pt]
\item Hern\'an, M. A. and Robins, J. M. (2023+). \textit{Causal Inference: What If}.
\item Huang, B., Zhang, K., Zhang, J., Glymour, C., and Sch\"olkopf, B. (2020). Causal discovery from heterogeneous and nonstationary data with independent changes.
\item Pearl, J. (2009). \textit{Causality}.
\item Rosenbaum, P. R. (2002). \textit{Observational Studies}.
\item Rosenbaum, P. R. (2020). \textit{Design of Observational Studies}.
\item Runge, J. et al. (2019). Detecting and quantifying causal associations in large nonlinear time series datasets.
\item Sadeghi, A., Gopal, A., and Fesanghary, M. (2024). \textit{Causal Discovery in Financial Markets: A Framework for Nonstationary Time-Series Data}.
\item VanderWeele, T. and Shpitser, I. (2011). A new criterion for confounder selection.
\end{itemize}

\end{document}
